\worksheettitle
  {Домашняя работа: ответы}
  {\modulemark{M07-H} Условия между цифрами}

\begin{teacherbox}[Критерии проверки]
  В заданиях 2--5 оцениваются отдельно: математическая запись условия,
  осознанный выбор перебираемой цифры, полный список и обоснование полноты.
  Случайно полученный правильный список без упорядоченного перебора не
  подтверждает достижение результата M07.
\end{teacherbox}

\section{Ответы}

\begin{enumerate}
  \item Записи условий:
    \begin{enumerate}
      \item \(b=a+3\);
      \item \(a=b+2\), или \(b=a-2\);
      \item \(b=2a\);
      \item \(a+b=11\).
    \end{enumerate}
  \item \(b=a+5\): перебираем \(a=1,\ldots,9\); \(16,27,38,49\).
    Кандидаты для \(a=5,\ldots,9\) больше \(9\).
  \item \(b=a-4\): можно перебирать \(a=1,\ldots,9\) или удобнее
    \(b=0,\ldots,5\); \(40,51,62,73,84,95\).
  \item \(a=3b\): перебираем \(b\); \(31,62,93\). При \(b=4\)
    получается \(a=12\), что не является цифрой.
  \item \(a+b=12\): перебираем \(a=1,\ldots,9\), сохраняем допустимые пары;
    \(39,48,57,66,75,84,93\). Всего 7 чисел.
  \item Ошибки:
    \begin{enumerate}
      \item \(10\) не цифра; полный список \(16,27,38,49\).
      \item Ноль допустим в единицах; полный список \(12,21,30\).
      \item Перепутаны цифры десятков и единиц; для \(a=3b\) верно
        \(31,62,93\).
    \end{enumerate}
  \item При сумме цифр \(9\) и числе меньше \(60\): \(18,27,36,45,54\).
  \item Условие: цифра единиц на \(4\) больше цифры десятков, \(b=a+4\).
    После \((5,9)\) получается пара \((6,10)\), а \(10\) не цифра.
\end{enumerate}

\section{Задание 9: оценивание объяснения}

Полное объяснение содержит четыре элемента:
\begin{enumerate}
  \item объяснён выбор цифры, изменяемой по порядку;
  \item перечислены все проверенные значения этой цифры;
  \item указано, как вычисляется вторая цифра;
  \item проверены ограничения на \(a\) и \(b\), отброшенные кандидаты объяснены.
\end{enumerate}

\newpage
\section{Решение о маршруте}

\begin{resultbox}[Проверка]
  \begin{itemize}
    \item \textbf{Результат устойчив:} задания 2--5 выполнены полными
      списками, направление зависимости не обращено, выбор перебираемой цифры
      осознан, объяснение содержит не менее трёх элементов.
    \item \textbf{Нужна M07-R:} повторяется обращение зависимости,
      используются значения вне диапазона цифр или пропускается число с нулём.
    \item \textbf{Нужна дополнительная практика:} условия записаны верно, но
      отдельные списки неполны или не обоснованы.
  \end{itemize}
\end{resultbox}
